\section{Introduction}

In 1955, Paul Erd\H{o}s \cite{someremarksonnumbertheory} posed the problem of estimating the number of distinct entries in an $N$-by-$N$ multiplication table. 
More precisely, letting 
\begin{equation*}
    M(N) = \#\set{ab}{1 \le a, b \le N}
\end{equation*}
denote the number of integers that can be written as the product of two integers up to $N$, what can be said about the growth rate of $M(N)$ as $N \rightarrow \infty$?
The statement of the problem is deceptively simple; yet it is not known whether there exists an asymptotic formula for $M(N)$.

\begin{figure}[h]
    \centering
    \begin{tabular}{c|cccccccc}
        & \textcolor{orange}{1} & \textcolor{orange}{2} & \textcolor{orange}{3} & \textcolor{orange}{4} & \textcolor{orange}{5} & \textcolor{orange}{6} & \textcolor{orange}{7} & \textcolor{orange}{\dots} \\ \hline
        \textcolor{orange}{1} & 1 & 2 & 3 & 4 & 5 & 6 & 7 & \dots \\ 
        \textcolor{orange}{2} & 2 & 4 & 6 & 8 & 10 & 12 & 14 & \dots \\ 
        \textcolor{orange}{3} & 3 & 6 & 9 & 12 & 15 & 18 & 21 & \dots \\ 
        \textcolor{orange}{4} & 4 & 8 & 12 & 16 & 20 & 24 & 28 & \dots \\ 
        \textcolor{orange}{5} & 5 & 10 & 15 & 20 & 25 & 30 & 35 & \dots \\ 
        \textcolor{orange}{6} & 6 & 12 & 18 & 24 & 30 & 36 & 42 & \dots \\ 
        \textcolor{orange}{7} & 7 & 14 & 21 & 28 & 35 & 42 & 49 & \dots \\ 
        \textcolor{orange}{$\vdots$} & $\vdots$ & $\vdots$ & $\vdots$ & $\vdots$ & $\vdots$ & $\vdots$ & $\vdots$ & $\ddots$ \\
    \end{tabular}
    \caption{The first entries of the multiplication table}
    \label{table:multiplication-table}
\end{figure}

Trivially, $N \le M(N) \le N(N + 1)/2$. 
Paul Erd\H{o}s showed that $M(N) = N^{2 - o(1)}$, see \cite{someremarksonnumbertheory} (in Hebrew) and \cite{anasymptoticinequality} (in Russian).
Further progress was made by Gérald Tenenbaum \cite{tenenbaumsurlaprobabilite} (in French).
However, the problem of determining the correct order of magnitude remained open for more than 50 years and was only solved in 2008 by Kevin Ford \cite{thedistributionofintegers, integerswithadivisor}.
He established the following theorem.

\begin{theorem}[{\cite[Corollary 3]{thedistributionofintegers}}] \label{thm:multiplication-table-problem}
    We have
    \begin{equation*}
        M(N) \asymp \frac{N^2}{(\log N)^\delta (\log \log N)^{3/2}},
    \end{equation*}
    where 
    \begin{equation*}
        \delta = 1 - \frac{1 + \log \log 2}{\log 2} \approx 0.086071.
    \end{equation*}
\end{theorem}

The goal of this dissertation is to give an exposition of Ford's work leading to the above result. 
In Chapter 2, we start by developing some general theory of the factorisation properties of ``typical'' integers that is relevant to the problem. 
This includes their number of prime divisors and the distribution thereof. 
Furthermore, we present some simple estimates for the multiplication table problem dating back to Erd\H{o}s.
After that, we give a detailed account of the proof of Theorem \ref{thm:multiplication-table-problem} in Chapter 3.
The main device to study $M(N)$ is the quantity $H(x, y, 2y)$, the number of integers up to $x$ with a divisor in $(y, 2y]$.
Ford's work reveals a deep connection between $H(x, y, 2y)$ and the distribution of prime divisors on a log-log scale, and much of Chapter 3 is concerned with studying that connection.

\paragraph{Notation.}
We denote by $\mN$ the set of natural numbers $\Set{1, 2, \dots}$, and by $[N]$ the set of integers $\Set{1, \dots, N}$.
The letter $p$ (as well as $p_i, p'$, and other variations) always denotes a prime number. 
% For example, $\sum_{p \le y}$ means that the summation is over all primes up to $y$.
The notation $p^\nu\parallel n$, for $\nu \in \mN$, means that $p^\nu \mid n$, but $p^{\nu + 1} \nmid n$. 
Letters in boldface (e.g. $\mbx, \mby$) always denote vector quantities.

The number of primes up to $x$ is denoted by $\pi(x)$.
For any natural number $n$, $P^-(n)$ and $P^+(n)$ denote the smallest, resp. largest, prime factor of $n$. 
We use the notational conventions $P^-(1) = \infty$ and $P^+(1) = 0$.
The number of distinct prime factors of $n$ is denoted by $\omega(n)$, and the number of divisors by $\tau(n)$. 
% Similarly, $\tau(n, y, z)$ is the number of divisors of $n$ in the interval $(y, z]$.

All logarithms are taken to the base $e$.
The notation $\log_k x$ is used to mean a $k$-fold iterated logarithm rather than a logarithm to the base $k$.
We will make frequent use of standard asymptotic notation as introduced in C3.8 Analytic Number Theory.

\paragraph{Preliminaries.}
We record a few standard results that will be used throughout the dissertation, starting with Stirling's formula:
\begin{equation*}
    n! \sim \sqrt{2 \pi n} \left( \frac{n}{e} \right)^n \quad \text{ as } n \rightarrow \infty.
\end{equation*}
A proof can be found in \cite[Pages 52--53]{feller}, for example.

We will make repeated use of the prime number theorem (PNT) in the form 
\begin{equation*}
    \pi(x) \sim \frac{x}{\log x} \quad \text{ as } x \rightarrow \infty.
\end{equation*}
We will also need a version of Mertens' first and second estimates with strong error terms. 
These can be deduced from the PNT with its classical de la Vallée-Poussin error term (as stated and proved in C3.8 Analytic Number Theory).

\begin{theorem} \label{thm:mertens-with-strong-error-terms}
    We have
    \begin{equation} \label{eq:mertens-first}
        \sum \limits_{p \le x} \frac{\log p}{p} = \log x + c_1 + O (e^{-c_0 \sqrt{\log x}})
    \end{equation}
    and 
    \begin{equation} \label{eq:mertens-second}
        \sum \limits_{p \le x} \frac{1}{p} = \log \log x + c_3 + O (e^{-c_2 \sqrt{\log x}}),
    \end{equation}
    for some absolute constants $c_0, c_1, c_2, c_3$.
\end{theorem}
As for Mertens' third estimate, the simple version 
\begin{equation*}
    \prod \limits_{p \le x} \left(1 - \frac{1}{p} \right) \asymp \frac{1}{\log x}
\end{equation*}
will be sufficient for our purposes.

Lastly, we will use the following general estimate, which is due to Halberstam and Richert (cf. \cite[Theorem 0.1]{divisors}).
\begin{theorem} \label{thm:halberstam-richert}
    Let $f: \mN \rightarrow \mR$ be a nonnegative multiplicative function.
    Suppose that there exist some $A$ and $B$ such that, for $x \ge 1$,
    \begin{align*}
        \sum \limits_{p \le x} f(p) \log p &\le Ax
    \end{align*}
    and
    \begin{align*}
        \sum \limits_{p} \sum \limits_{\nu \ge 2} \frac{f(p^\nu)}{p^\nu} \log (p^\nu) &\le B.
    \end{align*}
    Then, for all $x > 1$ one has
    \begin{equation*}
        \sum \limits_{n \le x} f(n) \ll_{A, B} \frac{x}{\log x} \sum \limits_{n \le x} \frac{f(n)}{n}.
    \end{equation*}
\end{theorem}
A proof is sketched in \cite[Exercise 00]{divisors}.